\chapter{LITERATURE REVIEW}

\section{Introduction}
Automated demographic occupancy counting in hospitals requires three fields of study: computer vision object detection, developmental allometry, and multi-object tracking. This chapter reviews the literature across these three areas. It evaluates how vision models handle human detection under heavy crowd density and occlusion. It analyzes biological growth rules that govern human body proportions. It reviews multi-target tracking algorithms. Finally, it identifies the research gaps that motivate this thesis.

\section{Evolution of Object Detection in Computer Vision}
Object detection evolved through two main phases: hand-crafted feature extractors and deep convolutional networks.

\subsection{Classical Hand-Crafted Feature Extractors}
Before deep learning, object detectors used hand-crafted spatial descriptors and shallow classifiers. Viola and Jones \cite{viola2001} built a fast face detector using Integral Images and Haar-like rectangular filters. By computing sum-area tables in constant time $\mathcal{O}(1)$, the Viola-Jones detector evaluated weak classifiers with AdaBoost in a cascade. This rejected non-face background areas quickly.

Later methods used gradient statistics, such as the Histogram of Oriented Gradients (HOG) and the Scale-Invariant Feature Transform (SIFT). In HOG, local gradient magnitudes and orientations form 1D histograms over pixel cells, normalized across overlapping blocks. Combined with linear Support Vector Machines (SVM), HOG set the baseline for pedestrian detection. However, hand-crafted descriptors showed clear weaknesses:
\begin{enumerate}[(i)]
    \item They used rigid spatial templates that failed under bodily deformations, seated postures, and angled carrying positions.
    \item They degraded under changing lighting, shadows, and low contrast in hospital hallways.
    \item Their computation scaled poorly with multi-scale sliding windows, preventing real-time multi-class detection on low-power CPUs.
\end{enumerate}

\subsection{Two-Stage Deep Detectors: Region-Based CNNs}
Deep Convolutional Neural Networks (CNNs), starting with AlexNet and formalized for detection in R-CNN, improved feature representation. Deep Residual Learning (ResNet) \cite{he2016resnet} introduced identity shortcut connections to solve vanishing gradients. Two-stage detectors divide object detection into two steps:
\begin{enumerate}[(i)]
    \item \textbf{Region Proposal Generation:} Finding class-agnostic regions of interest (RoIs) using Selective Search or Region Proposal Networks (RPN) in Faster R-CNN.
    \item \textbf{Classification and Bounding Box Refinement:} Extracting feature vectors via RoI Pooling or RoI Align and feeding them into fully connected layers for classification and coordinate regression.
\end{enumerate}

Two-stage detectors achieve high Mean Average Precision (mAP). However, processing candidate regions in two stages slows inference to 5--10 frames per second. This makes them too slow for continuous CCTV surveillance on edge hardware.

\subsection{One-Stage Deep Detectors: The YOLO Family}
To increase throughput, Redmon et al.\ \cite{redmon2016yolo} introduced the ``You Only Look Once'' (YOLO) framework. YOLO treats object detection as a single spatial regression task. Instead of proposing regions, YOLO divides the input image into an $S \times S$ grid. If an object center falls into a grid cell, that cell predicts bounding coordinates $(x, y, w, h)$, confidence $C$, and class probabilities $P(\text{Class}_i \mid \text{Object})$.

Later architectures, including modern anchor-free models (such as YOLOv8 and YOLO26s), added Path Aggregation Networks (PANet), Cross-Stage Partial (CSP) connections, and Task-Aligned Assigners \cite{jiang2022, jocher2023ultralytics}. These features fuse shallow high-resolution layers with deep semantic layers.

However, standard detectors trained on benchmarks like MS COCO group all humans into one ``person'' class. When fine-tuned on pediatric datasets such as the Child Detection Dataset (CDD) \cite{diop2025}, these models struggle with close physical proximity. Recent medical vision models have explored infant monitoring and spasm detection in specialized clinical wards \cite{diop2024}. Yet they often merge an infant carried on a mother's back into the adult's bounding box.

\section{Pediatric Anthropometry and Developmental Allometry}
Early computer vision studies estimated age groups from facial geometry \cite{kwon1999, liu2020}. However, facial features are frequently invisible in downward-tilted ceiling cameras. To separate children from adults reliably, vision models can evaluate whole-body allometric growth rules from physical anthropology \cite{bogin1999}.

\subsection{The Cephalocaudal Growth Gradient}
Human physical growth does not follow uniform isometric expansion. It follows allometric scaling rules \cite{bogin1999}. Allometry describes the relative growth rate of body parts compared to total body size:
\begin{equation}
    Y = a X^\alpha \implies \ln Y = \ln a + \alpha \ln X
\end{equation}
Here, $\alpha$ is the allometric growth coefficient. When $\alpha = 1$, growth is isometric. When $\alpha < 1$, growth is negative allometric. When $\alpha > 1$, growth is positive allometric.

Human growth follows the \textbf{cephalocaudal principle}, where growth progresses from head to feet \cite{bogin1999, mlinaric2024}:
\begin{itemize}
    \item \textbf{Cranial (Head) Development:} The skull and brain develop early. At birth, the infant head makes up approximately one-quarter ($25\%$) of total body height.
    \item \textbf{Somatic and Limb Development:} Postnatal growth centers on limb elongation, particularly the lower legs (femur, tibia, and fibula).
    \item \textbf{Adult Proportions:} By adulthood, limb growth reduces the head-to-body proportion to about one-eighth ($12.5\% - 13.3\%$) of total stature \cite{mlinaric2024}.
\end{itemize}

\subsection{Mathematical Implications for Computer Vision}
In 2D camera projections, absolute pixel height depends on camera distance and tilt angle. A child close to a ceiling camera can appear taller in pixels than an adult standing far away.

However, dimensionless body ratios remain stable under 2D perspective projections when the subject faces the camera plane:
\begin{equation}
    R_{head} = \frac{H_{head}}{H_{total}}, \qquad R_{ceph} = \frac{L_{torso}}{L_{leg}}
\end{equation}
17-point skeletal pose estimation \cite{cao2019openpose} extracts joint coordinates for the head, shoulders, hips, knees, and ankles. As Bogin \cite{bogin1999} and Mlinari\'{c} et al.\ \cite{mlinaric2024} show, pediatric subjects have $R_{head} \approx 0.25$ and $R_{ceph} \ge 0.85$. Adults have $R_{head} \approx 0.125$ and $R_{ceph} \le 0.70$. Adding these biological ratios directly into detection provides a clear mathematical rule to separate adults and children.

\section{Occlusion Handling in Crowded Clinical Healthcare Environments}
Physical occlusion is a major obstacle in visual monitoring. In outpatient clinics, waiting individuals sit close together on wooden benches or crowd around triage desks and pharmacy windows \cite{osei2024}.

\subsection{Taxonomy of Occlusion}
In hospital surveillance, occlusions fall into two groups:
\begin{enumerate}[(i)]
    \item \textbf{Static Environmental Occlusion:} Physical fixtures, including desks, registration counters, pillars, and carts, block occupant torsos and legs.
    \item \textbf{Dynamic Inter-Object Occlusion:} Walking or seated occupants block each other. In Ghana, mothers carry infants wrapped against their bodies \cite{agyepong1992}. This causes up to $80\%$ visual occlusion of the child \cite{ouyang2015}.
\end{enumerate}

\subsection{Algorithmic Strategies for Occlusion Robustness}
To handle heavy occlusion, modern vision research uses three methods:
\begin{itemize}
    \item \textbf{Intersection-over-Union (IoU) Loss Functions:} Standard bounding box regression used $\ell_1$ or $\ell_2$ coordinate errors, which do not reflect box overlap. Complete IoU (CIoU) \cite{zheng2020ciou} penalizes area overlap error, centroid distance, and aspect ratio differences. This stabilizes gradient updates under heavy overlap.
    \item \textbf{Slicing Aided Hyper Inference (SAHI):} In wide camera views, small children occupy small pixel patches (for example, $24 \times 24$ pixels). Downsampling an entire frame causes these small targets to disappear in pooling layers. Akyon et al.\ \cite{akyon2022sahi} introduced SAHI. SAHI cuts frames into overlapping high-resolution tiles, runs inference on each tile, and merges boxes with Non-Maximum Suppression. This preserves small targets.
    \item \textbf{Luminance Equalization with CLAHE:} Fluorescent lamps, dark skin tones, and dim corners reduce image contrast. Contrast Limited Adaptive Histogram Equalization (CLAHE) \cite{pizer1987clahe} enhances contrast in localized tiles and limits noise amplification along boundaries.
\end{itemize}

\section{Multi-Object Tracking and Temporal Association}
Single-frame detection alone cannot maintain an accurate count. A tracking system must connect detections across video frames to prevent counting the same person twice as occupants move through corridors.

\subsection{The Tracking-by-Detection Paradigm}
Multi-object tracking (MOT) typically uses the \textit{Tracking-by-Detection} approach. In this setup, a detector outputs bounding boxes in frame $t$. An association algorithm then matches these boxes to active tracks from earlier frames $1, \dots, t-1$.

\subsection{Evolution from SORT to ByteTrack}
Bewley et al.\ \cite{bewley2016sort} created Simple Online and Realtime Tracking (SORT). SORT uses a linear constant-velocity Kalman filter to predict positions and uses the Hungarian algorithm to match boxes by IoU distance. Beyond standard Kalman filtering, researchers have examined centroid-based tracking \cite{ralhan2023} and hierarchical graph structures \cite{koshkina2025} for identity persistence across occlusions. SORT is fast (over 200 FPS), but it drops lost tracks after only a few occluded frames.

Wojke et al.\ \cite{wojke2017deepsort} introduced DeepSORT. DeepSORT adds appearance embeddings from a deep Re-ID network to the matching cost. However, extracting embeddings for every box slows down CPU inference. Furthermore, when mothers and children wear identical patterned cloth, appearance embeddings become too similar, leading to ID switches.

Zhang et al.\ \cite{zhang2022bytetrack} introduced ByteTrack to solve this issue. Instead of dropping low-confidence boxes, ByteTrack matches in two stages:
\begin{enumerate}[(i)]
    \item It first matches high-confidence detections ($D_{high}$) to active tracks using spatial IoU and motion models.
    \item It then matches remaining tracks against low-confidence detections ($D_{low}$). This recovers occluded or blurred targets without creating false trajectories.
\end{enumerate}
ByteTrack tracks accurately using bounding coordinates alone, avoiding the latency of deep appearance embeddings.

\section{Edge Computing and Model Compression in Healthcare}
Deploying vision models in developing clinics introduces strict constraints. Facilities face unstable power, lack GPU servers, and prohibit transmitting raw video to cloud servers \cite{milkaite2019, addotey2023}.

\subsection{Knowledge Distillation Paradigms}
Knowledge distillation, developed by Hinton et al.\ \cite{hinton2015distilling}, transfers features from a large teacher model into a smaller student model. The loss balances hard ground-truth cross-entropy with softened teacher probabilities scaled by temperature $T$:
\begin{equation}
    \mathcal{L}_{KD} = (1 - \lambda) \mathcal{L}_{CE}(y, \sigma(z_s)) + \lambda T^2 \mathcal{L}_{KL}\left(\sigma\left(\frac{z_s}{T}\right), \sigma\left(\frac{z_t}{T}\right)\right)
\end{equation}
While distillation works well for image classification, applying it to dense object detection in specialized clinical domains presents challenges. As shown in Chapter 3, distillation quality depends directly on training data annotations. Noisy bounding boxes and occlusion in the student dataset cause the student network to copy and amplify errors from the teacher.

\subsection{Hardware SIMD Acceleration and ONNX Runtime}
Hospital edge devices require real-time CPU execution without discrete GPUs. Modern frameworks deploy models using ONNX Runtime \cite{onnxruntime2021} and Intel OpenVINO \cite{openvino2023} with SIMD acceleration.

Modern x86\_64 CPUs include Advanced Vector Extensions 2 (AVX2). AVX2 uses 256-bit registers that execute eight parallel 32-bit floating-point multiply-accumulate operations in one clock cycle:
\begin{equation}
    \text{AVX2 Throughput} = 8 \times \text{FP32 Multiply-Accumulate per Instruction}
\end{equation}
Lower precision formats like FP16 reduce memory use on GPUs. However, many edge CPUs lack native 16-bit float units and must emulate them in software, which slows inference down. Testing CPU SIMD execution profiles is therefore necessary for hospital deployment.

\section{Summary of Prior Art and What It Leaves Open}
Table \ref{tab:lit_comparison} summarizes the main detection and tracking approaches in the literature, comparing their core methods, strengths, and clinical limitations.

\begin{table}[htbp]
\centering
\caption{Comparison of Object Detection, Tracking, and Census Methods}
\label{tab:lit_comparison}
\resizebox{\textwidth}{!}{%
\begin{tabular}{lp{4.2cm}p{4.2cm}p{4.8cm}}
\toprule
\textbf{Method} & \textbf{Core Mechanism} & \textbf{Key Strengths} & \textbf{Clinical Limitations} \\
\midrule
Viola-Jones / HOG \cite{viola2001} & Haar wavelets / Gradient histograms + SVM & Low computational latency & Fails under non-rigid pose and cloth occlusion. \\
Faster R-CNN \cite{ouyang2015} & Two-stage RPN + RoI Align & High localization accuracy & Slow ($<8$ FPS), high memory use. \\
Standard YOLO \cite{diop2025, redmon2016yolo} & Single-stage anchor-free regression & Real-time inference speed & Merges carried infants into caregiver bounding boxes. \\
DeepSORT \cite{wojke2017deepsort} & Kalman filter + Deep Re-ID embeddings & Handles brief occlusions & Appearance features fail on identical cloth, high CPU load. \\
ByteTrack \cite{zhang2022bytetrack} & Two-stage Hungarian matching on $D_{high}, D_{low}$ & Recovers blurred targets without Re-ID & Lacks age awareness, requires secondary classification head. \\
\bottomrule
\end{tabular}%
}
\end{table}

\section{The Research Gap}
A review of the literature reveals three clear research gaps:
\begin{enumerate}[(i)]
    \item \textbf{The Semantic-Biological Disconnect:} Existing deep learning detectors optimize bounding boxes using geometric criteria (IoU, aspect ratio) alone. Reviewed literature does not integrate allometric cephalocaudal growth laws ($R_{\text{ceph}} = L_{\text{torso}}/L_{\text{leg}}$) directly as a verification constraint for pediatric person counting.
    \item \textbf{The Occlusion Feature-Collapse Gap:} Standard models run on downsampled full frames, which removes the fine visual details of carried infants. Among the literature evaluated, combining slicing-aided inference (SAHI) with local contrast enhancement (CLAHE) has not been widely examined for pediatric clinical triage.
    \item \textbf{The Data-Centric Distillation Dilemma:} While knowledge distillation offers compact models, its ability to transfer features for occluded children under clinical distribution shifts remains unproven.
\end{enumerate}

Addressing these three gaps forms the primary contribution of this thesis, detailed in Chapter 3.